\begin{table*}[t]
\centering
\scriptsize
\caption{Qualitative errors: base answer judged correct, best fine-tuned epoch incorrect}
\label{tab:T6_examples}
\begin{tabular}{p{1.6cm}p{1cm}p{2.4cm}p{2.6cm}p{3.2cm}p{3.2cm}p{2.4cm}}
\toprule
Model & Split & Question & Reference & Base answer & Fine-tuned answer & Contradicted key fact \\
\midrule
Qwen3-8B (ep1) & seen\_facts & At the macro level, what set of four overarching factors does the Unified Theory of Acceptance and Use of Technology put forward as shaping whether people take up a technology? & At the macro level, UTAUT identifies four higher-order determinants of technology adoption: performance expectancy, effort expectancy, social influence, and facilitating conditions. These synthesise earlier acceptance models, in contrast to the more granular cognitive emphasis of TAM2. & The Unified Theory of Acceptance and Use of Technology (UTAUT) identifies four overarching factors that shape whether people adopt and use a technology. These factors are: 1. **Performance Expectancy** – The degree to which an individual believes that using the technology will help them achieve … & The Unified Theory of Acceptance and Use of Technology posits that four macro-level factors shape technology adoption: cognitive appraisals (usefulness, ease of use), social influences (subjective norm, peer pressure), individual characteristics (self-efficacy, image), and emotional responses (satisfaction, trust). & UTAUT identifies four higher-order determinants of technology adoption: performance expectancy, effort expectancy, social influence, and facilitating conditions. \\
Qwen3-8B (ep1) & heldout\_docs & How does media framing of AI differ between the Global North and the Global South? & Global North media emphasises economic risks and efficiency concerns. Global South perspectives highlight potential societal benefits and democratisation opportunities. & Media framing of AI in the Global North often emphasizes innovation, economic competitiveness, and technological leadership, highlighting advancements in research, industry applications, and ethical governance. In contrast, media in the Global South tends to focus on AI's potential to address developmental challenges, such as healthcare, … & In the Global North, media coverage often highlights the transformative potential of AI, showcasing innovations in technology, business, and even creative industries. In contrast, in the Global South, media coverage tends to focus on the challenges of AI adoption, such as the digital divide, ethical … & Global North media emphasizes economic risks and efficiency concerns. \\
Qwen3-8B (ep1) & heldout\_docs & What is the approximate ratio of men to women in the tech industry, including ML model development, according to a widely cited fact in discussions of data diversity? & The ratio of men to women in tech is 70:30, and it has not changed much over the years. This lack of team diversity is a major challenge for the tech industry, including ML model development. & According to widely cited data, the ratio of men to women in the tech industry, including machine learning model development, is approximately **2:1**. This means men make up about **60-70\%** of the workforce, while women account for **30-40\%**. These figures are often referenced in discussions … & A widely cited fact states that only 2\% of the tech industry, including ML model development, is composed of women. & The ratio of men to women in the tech industry is 70:30. \\
Qwen3-8B (ep1) & seen\_facts & If a university in the UK is gauging how its master's-level business students feel about openness around AI, what did survey results from UK business postgraduates reveal about whether they … & In the survey, just over half of the participants, 51.1\%, identified a lack of transparency as an ethical concern. Almost half of the respondents did not consider it an ethical issue. & Survey results from UK business postgraduates have revealed that many consider opaque AI use to be an ethical issue. Students generally express concerns about transparency, accountability, and fairness in AI systems, particularly in contexts where decisions impact their professional practice or personal data. They often … & The survey found that 60\% of participants consider the use of AI with unclear levels of transparency as an ethical issue. This is a high percentage, and it shows that most participants consider the lack of transparency in AI an ethical issue. & 51.1\% of the participants identified a lack of transparency as an ethical concern. \\
Gemma 4 E4B (ep1) & heldout\_docs & What did a survey of education institutions find about how high schools differ from higher education institutions in developing generative AI policies? & Over 80\% of higher education institutions reported active policy development, 5\% already had a policy and 15\% had no plans, while only about 50\% of high schools were formulating policies and roughly 45\% had neither a policy nor plans. A … & A survey of education institutions generally found that **high schools are significantly less advanced and more hesitant than higher education institutions (HEIs) in developing comprehensive generative AI policies.** Key differences often highlighted include: * **Maturity of Policy:** HEIs tend to have more established, nuanced, and … & High schools are more likely to have developed generative AI policies than higher education institutions, which is a positive step toward addressing the challenges and opportunities presented by generative AI. & Over 80\% of higher education institutions reported active policy development. \\
Gemma 4 E4B (ep1) & heldout\_docs & Which organization introduced model cards, and in what year? & Model cards were introduced in 2019 by Google Research. & The **Model Cards for Model Reporting** were introduced by **Google** in **2019**. & The Institute for Ethical AI introduced model cards in 2019. & Google Research introduced model cards. \\
Gemma 4 E4B (ep1) & indomain & Which earlier concept did the TPACK framework grow out of, and who added Technological Knowledge to it? & The TPACK framework is rooted in Shulman's concept of Pedagogical Content Knowledge, which emphasizes integrating pedagogy and content knowledge to enhance student learning. Mishra and Koehler extended it by adding Technological Knowledge, resulting in TPACK with TCK, TPK, and the … & The TPACK framework grew out of **Mishra and Koehler's original framework**, which was an extension of **Shulman's Pedagogical Content Knowledge (PCK)**. **Mishra and Koehler** added the **Technological Knowledge (TK)** component to the existing knowledge structures to create the full TPACK model. & TPACK evolved from the original Technological Pedagogical Knowledge (TPK) framework, introduced by Berg and Scherer (2009). Technological Knowledge (TK) was added by Koehler and Mishra (2009). & The TPACK framework grew out of Shulman's concept of Pedagogical Content Knowledge. \\
Gemma 4 E4B (ep1) & seen\_facts & When screening for text produced by generative AI, what is meant by the burstiness metric, and what is the reason it can expose writing that a machine generated? & Burstiness, coupled with perplexity, is a measure of how frequently certain words or phrases appear in a text. It can identify language-use patterns indicative of machine-generated text, because GenAI tends to produce text that is less varied and more repetitive … & \#\# Burstiness Metric in AI Text Detection **What it means:** Burstiness, in the context of text analysis, refers to the **variation in sentence structure, length, and complexity** within a piece of writing. A high burstiness score indicates a significant fluctuation—a mix of very short, punchy … & Burstiness is the variation in sentence length and structure. A high burstiness score can indicate that the text was generated by a machine, because AI-generated text often has a more varied and complex structure than human-written text. & GenAI tends to produce text that is less varied and more repetitive than human writing. \\
Llama 3.1 8B (ep2) & heldout\_docs & What are the two components of Microsoft's open-source Fairlearn toolkit? & Fairlearn's two components are an interactive dashboard and unfairness mitigation algorithms. & The two components of Microsoft's open-source Fairlearn toolkit are: 1. **Auditors**: These are tools that help identify and measure bias in machine learning models. They provide metrics and insights to understand the fairness of a model. 2. **Mitigators**: These are algorithms and techniques that can … & Fairlearn has two components: metrics for assessing when a model is unfair, and metrics for comparing multiple models. & Fairlearn's two components are an interactive dashboard and unfairness mitigation algorithms. \\
Llama 3.1 8B (ep2) & indomain & Which earlier concept did the TPACK framework grow out of, and who added Technological Knowledge to it? & The TPACK framework is rooted in Shulman's concept of Pedagogical Content Knowledge, which emphasizes integrating pedagogy and content knowledge to enhance student learning. Mishra and Koehler extended it by adding Technological Knowledge, resulting in TPACK with TCK, TPK, and the … & The TPACK (Technological Pedagogical Content Knowledge) framework grew out of the PCK (Pedagogical Content Knowledge) framework. The PCK framework was developed by Lee Shulman in 1986. TPACK was later developed by Punya Mishra and Matthew J. Koehler in 2006. They added Technological Knowledge to the … & TPACK evolved from the earlier BTPK model proposed by Mishra and Koehler in 2006. Koehler and Mishra expanded BTPK to TPACK in 2007 by incorporating Technological Knowledge. & The TPACK framework grew out of Shulman's concept of Pedagogical Content Knowledge. \\
Llama 3.1 8B (ep2) & seen\_facts & After Congress passed the 1958 National Defense Education Act, by how much did the number of students enrolled in American colleges grow over the decade from 1960 to 1970? & In 1960, 3.6 million students attended college, and by 1970 the number was 7.5 million. Many students were able to attend college because of the National Defense Education Act, which authorized more than \$1 billion toward decreasing student loans, funding … & The number of students enrolled in American colleges grew by 96\% over the decade from 1960 to 1970. & Following the National Defense Education Act of 1958, the number of students enrolled in American colleges and universities increased by 60\% between 1960 and 1970, from 3.6 million to 5.7 million. & By 1970, 7.5 million students attended college. \\
Llama 3.1 8B (ep2) & heldout\_docs & What does EDIB stand for, and what must evaluation of LLM-generated educational text assess in relation to it? & EDIB stands for Equity, Diversity, Inclusion, \& Belonging. Evaluation of generated text must assess its EDIB appropriateness, along with its factual accuracy. & EDIB stands for Equity, Diversity, Inclusion, and Belonging. When evaluating LLM (Large Language Model)-generated educational text, it's essential to assess its alignment with EDIB principles. This involves evaluating the text for potential biases, stereotypes, and culturally insensitive content. The assessment should consider the following aspects: … & EDIB stands for Evaluation of Data Integrity, Bias. It involves examining the AI-generated content for accuracy, fairness, and any inherent biases that may be present. & EDIB stands for Equity, Diversity, Inclusion, \& Belonging. \\
\bottomrule
\end{tabular}
\par\smallskip\footnotesize Answers truncated to about 45 words. Contradicted fact = reference key fact with the highest NLI contradiction probability ($\geq$ 0.5) against the fine-tuned answer; – if none.
\end{table*}
